% Advanced Functional Materials (AFM) Template
% Publisher: Wiley-VCH
% Submission: ScholarOne Manuscripts portal
%
% Article types:
%   Full Paper  : ≤10 published pages (~7,500–8,000 words + figures)
%   Communication: ≤5 published pages (~3,500–4,000 words + figures)
%   Review      : ≤12,000 words
%
% Abstract (Full Paper): ≤200 words
% Abstract (Communication): opening paragraph acts as abstract (no separate section)
% Keywords: exactly 5
% Graphical abstract: required; separate file, 5 × 12.7 cm, 300 dpi
% Figures: TIFF/EPS/PDF, ≥300 dpi, with line numbers + double-spacing for review
% References: numbered style [1]

\documentclass[12pt]{article}

\usepackage[margin=1in]{geometry}
\usepackage{amsmath,amssymb,bm}
\usepackage{graphicx}
\usepackage{xcolor}
\usepackage{hyperref}
\usepackage{siunitx}
\usepackage{booktabs}
\usepackage{natbib}
\usepackage{authblk}
\usepackage[version=4]{mhchem}
\usepackage{lineno}    % required: line numbers throughout
\usepackage{setspace}  % required: double spacing
\doublespacing
\linenumbers

% Wiley/AFM prefers submission as a single compiled PDF with line numbers.
% Figures embedded in text for review. Provide separate high-res files at acceptance.

% Placeholder macro
\newcommand{\placeholder}[1]{%
  \begin{center}
    \fbox{\parbox{0.9\textwidth}{\centering\small\textit{[FIGURE: #1]}}}
  \end{center}
}

% ─── Title ───────────────────────────────────────────────────────────────────
% Must clearly state the functional advance; quantify performance if possible
\title{\textbf{[Functional advance verb]-Enabled [Material/Device] with [Quantified Performance] for [Application]}}

\author[1]{Author One}
\author[1,2]{Author Two}
\author[2*]{Corresponding Author}

\affil[1]{Department Name, University Name, City, Country}
\affil[2]{Institute Name, City, Country}
\affil[*]{E-mail: corresponding@email.com}

\date{}

\begin{document}

\maketitle

% ─── Keywords ────────────────────────────────────────────────────────────────
% Exactly 5 keywords required
\noindent\textbf{Keywords:} functional materials, keyword2, keyword3, keyword4, keyword5

% ─── Abstract (Full Paper) ───────────────────────────────────────────────────
% ≤200 words; emphasize the functional advance and measured performance gain
% For Communication: delete this block; opening paragraph IS the abstract
\begin{abstract}
Context sentence on the application need or materials challenge.
Problem sentence: what functional property is limited by current approaches.
Approach sentence: synthesis/fabrication/device strategy.
Key quantitative result: [X-fold improvement / record value] in [functional metric].
Broader significance: potential impact on [energy / electronics / biomedicine / etc.].
\end{abstract}

% ─── Introduction ────────────────────────────────────────────────────────────
\section{Introduction}

Broad functional materials context and application motivation.\cite{ref1}
Current state of the art and performance benchmarks.\cite{ref2,ref3}
Specific functional limitation or gap.\cite{ref4,ref5}
How this work addresses the gap: mechanism, synthesis strategy, or device concept.
``In this work, we demonstrate...'' with a preview of the key result.
Brief outline of the paper.

% ─── Results and Discussion ──────────────────────────────────────────────────
\section{Results and Discussion}

% AFM often combines Results and Discussion in a single section

\subsection{Synthesis / Fabrication and Structural Characterization}

\begin{figure}[htbp]
  \centering
  \placeholder{Synthesis scheme, XRD/SEM/TEM/AFM characterization panels}
  \caption{\textbf{Figure 1.} \textbf{Structural characterization.}
           a) Schematic illustration of the synthesis/fabrication route.
           b) XRD pattern confirming phase purity.
           c) SEM/TEM image (scale bar: \SI{200}{\nano\meter}).
           d) AFM topography and height profile.}
  \label{fig:structure}
\end{figure}

\subsection{Functional Properties and Performance}

\begin{figure*}[htbp]
  \centering
  \placeholder{Core functional performance: electrical, optical, catalytic, or mechanical data}
  \caption{\textbf{Figure 2.} \textbf{[Functional property] characterization.}
           a)--b) Measured [property] as a function of [parameter].
           c) Comparison with state-of-the-art benchmarks from the literature.
           d) Stability/cycling/durability data.}
  \label{fig:performance}
\end{figure*}

\subsection{Mechanism and Structure–Property Relationship}

\begin{figure}[htbp]
  \centering
  \placeholder{DFT/MD simulations or spectroscopic evidence for mechanism}
  \caption{\textbf{Figure 3.} \textbf{Mechanistic understanding.}
           a) DFT band structure / DOS / charge density difference.
           b) XPS / Raman / NEXAFS spectra comparing pristine and functional state.
           c) Proposed mechanism schematic.}
  \label{fig:mechanism}
\end{figure}

\subsection{Device Demonstration or Application}

\begin{figure}[htbp]
  \centering
  \placeholder{Device performance metrics and benchmark comparison}
  \caption{\textbf{Figure 4.} \textbf{Device performance.}
           a)-d) [Device metrics: efficiency, sensitivity, power density, etc.]}
  \label{fig:device}
\end{figure}

% ─── Conclusion ──────────────────────────────────────────────────────────────
\section{Conclusion}

In summary, we have demonstrated [functional advance] by [mechanism/strategy].
Key achievement: [quantified performance metric compared to prior benchmark].
The structure–property relationship reveals [insight].
This work opens a path toward [practical application].

% ─── Experimental Section ────────────────────────────────────────────────────
\section{Experimental Section}
% AFM: Experimental Section after Conclusion

\subsection{Materials}

\subsection{Synthesis Procedure}

\subsection{Characterization}
% List all instruments with model numbers and settings

\subsection{Device Fabrication}

\subsection{Electrochemical / Optical / Mechanical Measurements}

\subsection{Computational Details}
% Software, functional, basis set, convergence criteria

% ─── Supporting Information ──────────────────────────────────────────────────
% Separate file; referenced as "Supporting Information Figure S1"
\noindent\textbf{Supporting Information}

Supporting Information is available from Wiley Online Library or from the author.

% ─── Acknowledgements ────────────────────────────────────────────────────────
\section*{Acknowledgements}
The authors acknowledge...
This work was supported by [Agency] (Grant No. XXX).

\noindent\textbf{Author Contributions} (CRediT):
A.O.: conceptualization, methodology, writing—original draft.
C.A.: supervision, funding acquisition, writing—review \& editing.

\noindent\textbf{Conflict of Interest:} The authors declare no conflict of interest.

\noindent\textbf{Data Availability Statement:} Research data are available from the corresponding author on reasonable request.

% ─── References ──────────────────────────────────────────────────────────────
% AFM: numbered style [1]; use natbib unsrt or similar
\bibliographystyle{unsrt}
\bibliography{refs}

\end{document}
